%═══════════════════════════════════════════════════════════════════════════════
\section{Theoretical Framework}\label{sec:decision}
%═══════════════════════════════════════════════════════════════════════════════

\subsection{The Decision Problem}\label{sec:problem}

Let $Z$ be an unknown outcome in a finite outcome space $\mathcal Z$, and let
$d\in\mathcal D$ be an action chosen before $Z$ is observed. The loss
$L(d,z)\geq0$ is what choosing $d$ costs when the outcome turns out to be $z$.
The decision-maker supplies this loss table, and it records which mistakes
matter and by how much. With finitely many outcomes the table is a grid of
numbers, one per action and outcome, and $L(d,\cdot)$ denotes the row for
action $d$, listing what that action costs in each outcome in turn.

Uncertainty about outcome $Z$ enters in two separate places. The first is the model $\mathcal P$, a set of distributions any one of which
could be the true one. The decision-maker is unwilling to stake the guarantee
on a single distribution, so the report is required to be valid under every
member at once. As $\mathcal P$ gets larger, the statement about what
the world might be doing becomes more cautious, and it buys a weaker guarantee in return.

The second is the belief $\widehat P_1$, the single distribution the
decision-maker actually thinks is true, formed before any data is collected.
It carries no guarantee whatever. Intuitively, it is simply what they would act on ``if
nothing could go wrong''.

The two do different jobs. The model says what the report must survive, the
belief says what it should be good at. They are not competing estimates of one
quantity, and the thesis is about what happens when they disagree about which
action to take. Two standing assumptions apply throughout: every distribution
in $\mathcal P$ gives every outcome positive probability, so each expected
loss is a finite sum, and every expectation is taken under the distribution
named in its subscript.

\subsection{Evidence Profiles and Validity}\label{sec:evar}

A profile $\mathcal E$ is a valid e-variable for the model $\mathcal P$ when
\begin{equation}
  \mathbb E_{P}[\mathcal E(Z)]\leq1
  \qquad\text{for every }P\in\mathcal P .
  \label{eq:evalue_def}
\end{equation}

Read as a budget, the condition gives the profile one unit of e-value to spend
on average. Nothing dictates how that unit is spread across outcomes, and how to spread it is precisely the design problem of this
thesis. Constructing a valid profile from raw data is a separate problem,
solved by the methods cited in Section~\ref{sec:introduction}, and is taken as
given here.

What makes the condition worth accepting is what it buys. Markov's inequality
turns it into
\begin{equation}
  P[\mathcal E(Z)\geq1/\alpha]\leq\alpha
  \qquad\text{for every }\alpha\in(0,1),
  \label{eq:markov}
\end{equation}
so a profile whose score at an outcome exceeds $1/\alpha$ can be treated as
evidence at level $\alpha$ against it. This holds at every $\alpha$
simultaneously and under every member of the model, so a single profile serves
every threshold a reader might apply with no correction for inspecting
several. A companion result \citep{Grunwald2024,Koning2024posthoc} extends it
to thresholds chosen after the scores are seen, so evidence can be published
before anyone fixes how strictly to read it.

One further consequence matters below. Only the hardest constraint in
\eqref{eq:evalue_def} does any work: the profile is valid exactly when its
largest expectation across the model is at most one, so validity is settled by
the single member under which the profile looks largest. A profile
that spends the budget exactly under some member, rather than leaving slack
everywhere, is the interesting case, since any leftover could be spent on
sharper evidence. Every profile constructed in this thesis is of that kind. The
examples below use the simplest model, a single distribution
$\mathcal P=\{P_0\}$, to keep the arithmetic checkable, and
Section~\ref{sec:market} uses a genuinely composite one.

\subsection{From Prediction Sets to the Weighted Rule}\label{sec:rule}

A profile has to become an action. For a prediction set this is done by the
rule \eqref{eq:intro_set_rule}, which protects against the worst outcome the
set still admits. The profile analogue replaces the worst case over a set by a
worst case weighted by the evidence. This is the weighted worst-case rule of
\citet{KoningvanMeer2025},
\begin{equation}
  \delta(\mathcal E)\in\argmin_{d\in\mathcal D}C(d;\mathcal E),
  \qquad
  C(d;\mathcal E):=\max_{z\in\mathcal Z}
  \frac{L(d,z)}{\mathcal E(z)},
  \label{eq:wl_decision}
\end{equation}
with the conventions $x/0=\infty$ for $x>0$ and $0/0=0$. Dividing by the score
does not declare an outcome impossible: it discounts a loss that would be
incurred only if the outcome occurred, in proportion to the evidence against
it.

The number $C(d;\mathcal E)$ is called the certificate of $d$, because it
bounds expected loss. It depends only on the loss row and the profile, so it
is known before the outcome is. Assume $C(d;\mathcal E)<\infty$, since otherwise the
bound is vacuous. By definition of the maximum,
$L(d,z)\leq C(d;\mathcal E)\,\mathcal E(z)$ at every outcome, so for any
$P\in\mathcal P$,
\begin{equation}
  \mathbb E_{P}[L(d,Z)]
  \leq C(d;\mathcal E)\,\mathbb E_{P}[\mathcal E(Z)]
  \leq C(d;\mathcal E).
  \label{eq:cert_bound}
\end{equation}
The final inequality is where validity is used, and without it a profile could
be scaled up to make any certificate small. The bound holds under every
member of the model at once, and it concerns expected loss, not every realised
loss. The same pointwise inequality also gives a tail bound. A loss of at
least $C(d;\mathcal E)/\alpha$ forces a score of at least $1/\alpha$, so by
\eqref{eq:markov}
\[
  P\bigl[L(d,Z)\geq C(d;\mathcal E)/\alpha\bigr]\leq\alpha
  \qquad\text{for every }\alpha\in(0,1)\text{ and every }P\in\mathcal P.
\]

The rule may have several minimisers. An action is weakly selected when
it is one of them, and strictly selected when its certificate beats every
competitor. That distinction matters once an action is imposed as a constraint
in Section~\ref{sec:result2}.

\subsection{A Worked Example}\label{sec:setup}

A bettor stakes money on a football match with two outcomes, a win for $A$ or
a win for $B$. They may back $A$ through $d_1$ or hedge the position through
$d_2$. A losing bet costs twice as much as a hedge that turns out to have been
unnecessary. Listing each action's loss if $A$ wins and then if $B$ wins,
\begin{equation}
  L(d_1,\cdot)=(0,8),\qquad L(d_2,\cdot)=(4,0),
  \label{eq:two_loss}
\end{equation}
so backing $A$ loses $8$ only if $B$ wins, and the hedge loses $4$ only if $A$
wins. The model is the single distribution $\mathcal P=\{P_0\}$ with
$P_0=(1/2,1/2)$, an even match, and the bettor's belief is
$\widehat P_1=(0.75,0.25)$, favouring $A$.

The two criteria already disagree about what to do. The certificate bounds
expected loss under the model \eqref{eq:cert_bound}, so a decision-maker who
wants that guarantee compares actions by their average loss against $P_0$,
\[
\begin{aligned}
\mathbb E_{P_0}[L(d_1,Z)]&=\tfrac12(0)+\tfrac12(8)=4,\\
\mathbb E_{P_0}[L(d_2,Z)]&=\tfrac12(4)+\tfrac12(0)=2,
\end{aligned}
\]
and the hedge $d_2$ comes out ahead. A decision-maker who trusts the belief
averages under $\widehat P_1$ instead,
\[
\begin{aligned}
\mathbb E_{\widehat P_1}[L(d_1,Z)]&=0.75(0)+0.25(8)=2,\\
\mathbb E_{\widehat P_1}[L(d_2,Z)]&=0.75(4)+0.25(0)=3,
\end{aligned}
\]
so the belief favours backing $A$. Neither answer is wrong. They answer
different questions, and the decision-maker has to say which one they are
asking.

Now compare the two decision rules on this example. A profile becomes a
prediction set by keeping every outcome whose score stays below $1/\alpha$,
since by \eqref{eq:markov} a larger score counts as evidence at level $\alpha$
against it. Fix $\alpha=1/2$, so the threshold is $2$, and take two valid
profiles: the
flat profile $\mathcal E^{\mathrm{flat}}=(1,1)$, which spends its budget
evenly, and a graded profile $\mathcal E^{\mathrm{graded}}=(2/3,4/3)$, which
shifts weight onto $B$. Both are valid, since each averages to one under
$P_0$.

Neither profile reaches the threshold at either outcome, so both produce the
same prediction set $S=\{A,B\}$, and \eqref{eq:intro_set_rule} cannot tell them apart.
In both cases the worst loss over $S$ is $8$ for the bet and $4$ for the
hedge, so each backs the hedge.

The weighted rule separates them, because it divides outcome by outcome before
taking the worst case. Under the flat profile,
\[
  C(d_1)=\max\bigl(0/1,\; 8/1\bigr)=8,
  \qquad
  C(d_2)=\max\bigl(4/1,\; 0/1\bigr)=4,
\]
so $d_2$ wins outright, exactly as the prediction set concluded. The graded
profile changes the divisors,
\[
  C(d_1)=\max\bigl(0\big/\tfrac23,\; 8\big/\tfrac43\bigr)=6,
  \qquad
  C(d_2)=\max\bigl(4\big/\tfrac23,\; 0\big/\tfrac43\bigr)=6,
\]
and the two are now equal, so $d_1$ becomes available. The extra score on $B$
deflates the only outcome in which backing $A$ costs anything, while the score
removed from $A$ inflates the cost of hedging. The consequence is a different
action, not merely a different number, and under the belief $d_1$ costs $2$
where $d_2$ costs $3$. The dependence on $\alpha$ disappears too, since
\eqref{eq:wl_decision} never thresholds the profile, so no coverage level
enters the choice of action.

This is the whole problem in miniature: two valid profiles recommend
different actions, and validity does not choose between them.

A second example shows that the disagreement can also turn on a single
outcome that the belief considers rare but that is costly for most actions.
Showing this needs a third outcome, a draw, which settles both side bets as
losses at once. The bettor may back either side through $d_A$ or $d_B$, or
insure against the draw through $d_C$, which costs a premium whenever the
match has a winner and pays only in the draw. Writing the outcomes as a win
for $A$, a win for $B$, and a draw $C$,
\begin{equation}
  L(d_A,\cdot)=(0,2,10),\quad
  L(d_B,\cdot)=(2,0,10),\quad
  L(d_C,\cdot)=(4,4,0),
  \label{eq:three_loss}
\end{equation}
with $\mathcal P=\{P_0\}$ for $P_0=(1/3,1/3,1/3)$ and belief
$\widehat P_1=(0.6,0.3,0.1)$. The belief is confident about the side and
dismissive about the draw, putting $0.60$ on $A$ but only $0.10$ on $C$, which is
plausibly too low for football, where draws are common.

The two criteria disagree here too. Against the model the averages are
$(4,4,8/3)$, so protection points to $d_C$, while under the belief they are
$(1.6,2.2,3.6)$, so the belief points to $d_A$. The draw decides it: $d_A$ and $d_B$ each
expose the bettor to a loss of $10$ there, which $d_C$ avoids at the cost of never doing
well. Section~\ref{sec:sim_danger} asks how wrong such a belief can be before
trusting it stops paying.
